\section[Performance Metrics]{Performance Metrics and Design Trade-offs}
\label{sec:pd_metrics}

The transport and absorption physics of \cref{sec:pd_fundamentals} support a set
of engineering metrics that define what a useful photodetector must accomplish and
at what cost. These metrics are not independent: they form a multi-dimensional
trade-off space in which improving one quantity typically degrades another. Every
architectural and material choice in the sections that follow is justifiable
using the definitions established here.

\subsection[Responsivity and Quantum Efficiency]{Responsivity and Quantum Efficiency}
\label{subsec:pd_responsivity}

The \emph{external responsivity} $\mathcal{R}$ is:
\begin{equation}
    \mathcal{R} \equiv \frac{I_{\mathrm{ph}}}{P_{\mathrm{opt}}}
    \quad [\si{\ampere\per\watt}],
    \label{eq:responsivity}
\end{equation}
with theoretical upper bound:
\begin{equation}
    \mathcal{R}_{\max} = \frac{q\lambda}{hc}
    \approx \SI{1.057}{\ampere\per\watt}
    \quad\text{at \SI{1310}{\nano\metre}}.
    \label{eq:rmax}
\end{equation}
The external quantum efficiency (EQE) $\eta_{\mathrm{ext}}$ and internal quantum
efficiency (IQE) $\eta_{\mathrm{int}}$ are:
\begin{align}
    \eta_{\mathrm{ext}}
    &= \mathcal{R}\cdot\frac{hc}{q\lambda},
    \label{eq:eqe}\\
    \eta_{\mathrm{int}}
    &= \frac{I_{\mathrm{ph}}/q}{P_{\mathrm{abs}}/\hbar\omega}.
    \label{eq:iqe}
\end{align}
In the fully depleted intrinsic region at saturation field, the drift transit
time $\tau_{\mathrm{tr}}=W_i/v_{\mathrm{d,sat},h}$ is much shorter than the
SRH lifetime $\tau_0$, giving a carrier collection efficiency
$\xi\approx1-\tau_{\mathrm{tr}}/\tau_0\approx1$. The more general form of the
EQE chain is:
\begin{equation}
    \eta_{\mathrm{ext}}
    = (1 - R_{\mathrm{refl}})\,\xi\,\eta_c\,
      \bigl(1 - e^{-\alpha(\lambda)\,L}\bigr),
    \label{eq:qe_full}
\end{equation}
where $R_{\mathrm{refl}}$ is the Si/Ge entry-facet reflectance and $\eta_c$ is
the coupling efficiency from the silicon waveguide into the Ge absorber mode.
Real devices fall below $\mathcal{R}_{\max}$ because of coupling loss ($\eta_c<1$),
incomplete absorption ($e^{-\alpha L}>0$), and residual carrier recombination
before collection ($\xi<1$).

\begin{figure}[H]
    \centering
    \missingfigure[figwidth=0.84\linewidth]{
        Two-panel metrics figure. Left panel: photocurrent versus input
        optical power, showing the linear responsivity slope in the low-power
        regime and saturation at high power. Right panel: responsivity versus
        wavelength, highlighting the short-wavelength reduction from excess
        carrier energy and the long-wavelength roll-off near the absorption
        edge.}
    \caption{Responsivity viewed in both its power-domain and spectral forms.
    The low-power slope of the $I$-$P$ curve defines $\mathcal{R}$, whereas
    the wavelength dependence determines the useful operating window around the
    O-band.}
    \label{fig:responsivity_concept}
\end{figure}

\subsection[Bandwidth: Transit-Time and RC Limits]{Bandwidth: Transit-Time and RC Limits}
\label{subsec:pd_bw}

The frequency response is limited by two physically distinct mechanisms.
In the absence of deliberate inductive peaking, the combined baseline
bandwidth is~\cite{SalehTeich1991}:
\begin{equation}
    f_{3\,\mathrm{dB}}
    = \left(\frac{1}{f_{\mathrm{tt}}^2}
      + \frac{1}{f_{\mathrm{RC}}^2}\right)^{-1/2}.
    \label{eq:bw_total}
\end{equation}

\paragraph{Transit-time limit.}
With hole saturation velocity $v_{\mathrm{d,sat},h}\approx\SI{4e6}{\centi\metre\per\second}$
and electron saturation velocity $v_{\mathrm{d,sat},e}\approx\SI{6e6}{\centi\metre\per\second}$,
the transit-time-limited 3\,dB bandwidth from Fourier analysis of a rectangular
current impulse response is:
\begin{equation}
    f_{\mathrm{tt}} = \frac{0.45\,v_{\mathrm{d,sat},h}}{W_i},
    \label{eq:bw_tt}
\end{equation}
where the coefficient $0.45$ arises from the sinc-like spectral roll-off and
reflects the asymmetric electron-hole velocity ratio
$v_{\mathrm{d,sat},e}/v_{\mathrm{d,sat},h}\approx1.5$.
For $W_i=\SI{350}{\nano\metre}$: $f_{\mathrm{tt}}\approx\SI{51}{\giga\hertz}$.
This value is a geometry-screening lower bound rather than the absolute transport
ceiling. The full two-carrier response of \cref{eq:jt_analytical} is broader
because electrons and holes induce terminal current simultaneously, and the
experimentally relevant RF response further benefits from on-electrode inductance
$L_p$. For the \SI{350}{\nano\metre} intrinsic region, the benchmark-calibrated
transport ceiling is in the \SIrange{110}{120}{\giga\hertz} class.

\paragraph{RC limit.}
The junction capacitance of the fully depleted intrinsic region is:
\begin{equation}
    C_j = \frac{\varepsilon_0\varepsilon_{\mathrm{Ge}}\,A_j}{W_i},
    \label{eq:cj}
\end{equation}
where $\varepsilon_{\mathrm{Ge}}\approx16.2$ and $A_j=W_m\times L$.
The RC-limited bandwidth is:
\begin{equation}
    f_{\mathrm{RC}} = \frac{1}{2\pi R_{\mathrm{tot}}\,C_j},
    \label{eq:bw_rc}
\end{equation}
where $R_{\mathrm{tot}}=R_s+R_L$ and $R_L=\SI{50}{\ohm}$.

\paragraph{Optimal intrinsic-layer thickness.}
Setting $f_{\mathrm{tt}}(W_i^*)=f_{\mathrm{RC}}(W_i^*)$ and substituting
\cref{eq:bw_tt,eq:bw_rc,eq:cj} gives:
\begin{equation}
    W_i^* = \sqrt{0.9\pi\,R_{\mathrm{tot}}\,\varepsilon_0\varepsilon_{\mathrm{Ge}}
             \,v_{\mathrm{d,sat},h}\,A_j}.
    \label{eq:wi_optimal}
\end{equation}
Substituting $R_{\mathrm{tot}}=\SI{60}{\ohm}$, $\varepsilon_{\mathrm{Ge}}=16.2$,
$v_{\mathrm{d,sat},h}=\SI{4e6}{\centi\metre\per\second}$, and
$A_j=\SI{6}{\micro\metre\squared}$ gives
$W_i^*\approx\SI{310}{\nano\metre}$, within \SI{11}{\percent} of the chosen
$W_i=\SI{350}{\nano\metre}$, confirming that the device is positioned near the
bandwidth optimum. The expression \cref{eq:wi_optimal} reveals a fundamental
scaling law: since $W_i^*\propto\sqrt{A_j}$, shrinking the junction area by a
factor of four allows $W_i^*$ to be halved, raising the balanced bandwidth at
the new optimum by approximately the same factor.

The small-signal photocurrent transfer function is well-approximated by the
factorised form:
\begin{equation}
    H_{\mathrm{PD}}(f)
    = \mathrm{sinc}\!\left(\frac{f}{f_{\mathrm{tt}}}\right)
    \cdot\frac{1}{1+i\,f/f_{\mathrm{RC}}},
    \label{eq:H_pd}
\end{equation}
where $\mathrm{sinc}(x)=\sin(\pi x)/(\pi x)$. For the benchmark-calibrated
U-shaped device, the electrical sub-network is more accurately represented by a
series branch $Z_s(\omega)=R_s+j\omega L_p$ driving a terminated capacitive load
$Z_L(\omega)=\left(R_L^{-1}+j\omega C_p\right)^{-1}$ while the reverse-biased
junction impedance is $Z_j(\omega)\approx(j\omega C_j)^{-1}$. Current division
then gives the peaked electrical transimpedance:
\begin{equation}
    H_{\mathrm{elec,pk}}(\omega)
    = \frac{Z_j(\omega)\,Z_L(\omega)}
           {Z_j(\omega)+Z_s(\omega)+Z_L(\omega)},
    \label{eq:helec_pk}
\end{equation}
and the full benchmark-calibrated detector response becomes:
\begin{equation}
    H_{\mathrm{PD,pk}}(\omega)
    = H_{\mathrm{tt}}(\omega)\,H_{\mathrm{elec,pk}}(\omega).
    \label{eq:H_pd_pk}
\end{equation}
\Cref{eq:H_pd} is used for fast trade-off screening; \cref{eq:H_pd_pk} is used
when projecting the validated single-U device toward the \SI{100}{\giga\hertz}-class
operating regime.

\paragraph{Series resistance and RC time constant.}
The series resistance $R_s$ in the vertical n-i-p architecture is dominated by
the lateral conduction path through the $p^{++}$-Si platform:
\begin{equation}
    R_{s,\mathrm{anode}} \approx \frac{\rho_{\mathrm{Si}}}{t_{\mathrm{Si}}}
    \cdot\frac{\ell_a}{w_a},
    \label{eq:rs_anode}
\end{equation}
where $t_{\mathrm{Si}}=\SI{220}{\nano\metre}$ is the platform thickness.
The RC time constant is:
\begin{equation}
    \tau_{\mathrm{RC}} = R_s\,C_j,
    \label{eq:tau_rc}
\end{equation}
from which $f_{\mathrm{RC}}=(2\pi\tau_{\mathrm{RC}})^{-1}$.

\begin{figure}[H]
    \centering
    \missingfigure[figwidth=0.82\linewidth]{
        Two-panel RC figure. Left panel: simplified equivalent circuit showing
        $C_j$, $R_s$, parasitic capacitance, and the \SI{50}{\ohm} load.
        Right panel: qualitative capacitance and RC-limited bandwidth trends
        versus junction area, illustrating the penalty paid when optical
        absorption is pursued through a larger footprint.}
    \caption{Capacitance-resistance interpretation of the vertical PIN
    photodetector. The key RC penalty comes from simultaneously increasing the
    depleted junction area and the lateral access resistance, which is why
    layout and footprint must be optimised together.}
    \label{fig:rc_concept}
\end{figure}

\subsection[Noise, Sensitivity, and Figures of Merit]{Noise, Sensitivity, and Figures of Merit}
\label{subsec:pd_noise}

\subsubsection{Dark Current}
\label{subsubsec:pd_dark}

Dark current $I_d$ flows through the reverse-biased junction in the absence of
optical excitation and sets the noise floor through \cref{eq:shot}. Three
generation-recombination mechanisms contribute~\cite{SzeNg2006}. The \emph{SRH
generation} current density within the depleted intrinsic layer of thickness
$W_i$ follows from \cref{eq:srh_dep}:
\begin{equation}
    J_{\mathrm{SRH}} = q\,U_{\mathrm{SRH,dep}}\,W_i
    = \frac{q\,n_i\,W_i}{2\tau_0}.
    \label{eq:jdark_srh}
\end{equation}
\emph{Minority-carrier diffusion} from the quasi-neutral contact regions
contributes:
\begin{equation}
    J_{\mathrm{diff}}
    = \frac{qD_p\,n_i^2}{L_p\,N_D} + \frac{qD_n\,n_i^2}{L_n\,N_A},
    \label{eq:jdark_diff}
\end{equation}
where $L_p=\sqrt{D_p\tau_p}$ and $L_n=\sqrt{D_n\tau_n}$.
\emph{Interface recombination} at the Ge/Si heterojunction contributes:
\begin{equation}
    J_{\mathrm{surf}} \approx \frac{q\,n_i\,v_s}{2}.
    \label{eq:jdark_surf}
\end{equation}
The total dark current, including both bulk and surface (perimeter)
contributions, is:
\begin{equation}
    I_d = A_j\!\left(J_{\mathrm{SRH}} + J_{\mathrm{diff}}\right)
    + P_j\,q\,n_i\,v_s/2,
    \label{eq:idark_decomposed}
\end{equation}
where $P_j=2(W_m+L)$ is the junction mesa perimeter. This perimeter scaling has
a direct consequence for miniaturisation: as $A_j$ is reduced to decrease $C_j$
and improve bandwidth, the area-to-perimeter ratio $A_j/P_j$ decreases, increasing
the fractional contribution of the surface term.

\begin{figure}[H]
    \centering
    \missingfigure[figwidth=0.82\linewidth]{
        Two-panel dark-current figure. Left panel: decomposition of total dark
        current into SRH, diffusion, and surface terms versus reverse bias.
        Right panel: area scaling versus perimeter scaling for large and small
        junctions, showing why surface effects grow in importance as the mesa
        footprint shrinks.}
    \caption{Dark-current mechanisms of the reverse-biased Ge-on-Si PIN.
    The present device remains bulk-SRH dominated at the chosen footprint, but
    perimeter-driven surface generation becomes progressively more important as
    the junction is aggressively miniaturised.}
    \label{fig:dark_current_concept}
\end{figure}

\subsubsection{Noise Spectral Densities and Noise-Equivalent Power}
\label{subsubsec:pd_noise_psd}

\begin{figure}[H]
    \centering
    \missingfigure[figwidth=0.86\linewidth]{
        Four-panel noise concept figure. Discrete carrier arrivals, Poisson
        counting statistics, time-domain current fluctuations, and flat shot
        versus thermal noise spectral densities are shown on a common causal
        path from charge discreteness to electrical noise PSD.}
    \caption{Figure-first explanation of noise origin. Shot noise is not an
    abstract fitting parameter but the continuum limit of the counting
    fluctuations of discrete carriers arriving at the terminal.}
    \label{fig:shot_noise_origin}
\end{figure}

\paragraph{Origin of shot noise from carrier counting statistics.}
The physical origin of shot noise is the randomness of discrete carrier arrivals.
Let $N_T$ denote the number of carriers collected during an observation interval
$T$. If the arrivals are statistically independent, then $N_T$ is Poisson
distributed with mean $\mathbb{E}[N_T]=IT/q$ and variance
$\mathrm{Var}[N_T]=IT/q$. The collected charge is $Q_T=qN_T$, so:
\begin{equation}
    \mathrm{Var}[Q_T] = qIT.
    \label{eq:var_charge_poisson}
\end{equation}
Dividing by $T^2$ converts charge variance to variance of the time-averaged current:
\begin{equation}
    \mathrm{Var}\!\left[\bar{I}_T\right]
    = \frac{qI}{T}.
    \label{eq:var_current_poisson}
\end{equation}
Passing from the finite-time variance to the one-sided continuous spectral
density yields the standard result $S_{\mathrm{shot}}=2qI$.

\emph{Shot noise} arises from the quantum-mechanical discreteness of charge
carriers:
\begin{equation}
    S_{\mathrm{shot}}(f) = 2q\!\left(I_d + I_{\mathrm{ph}}\right)
    \quad [\si{\ampere\squared\per\hertz}].
    \label{eq:shot}
\end{equation}
\emph{Johnson (thermal) noise} from the load resistance contributes:
\begin{equation}
    S_{\mathrm{thermal}}(f) = \frac{4kT}{R_L}
    \quad [\si{\ampere\squared\per\hertz}].
    \label{eq:thermal}
\end{equation}
Integrating over the detection bandwidth gives the total photodetector noise
variance:
\begin{equation}
    \sigma_n^2 = i_n^2\,B_e
    \quad [\si{\ampere\squared}],
    \label{eq:sigma_n}
\end{equation}
where $i_n^2 = S_{\mathrm{shot}}$. The \emph{noise-equivalent power} (NEP) and
\emph{specific detectivity} $D^*$ are:
\begin{align}
    \mathrm{NEP}
    &= \frac{\sqrt{S_{\mathrm{shot}} + S_{\mathrm{thermal}}}}{\mathcal{R}}
    \quad [\si{\watt\per\sqrt\hertz}],
    \label{eq:nep}\\
    D^*_{\mathrm{shot}}
    &= \frac{\mathcal{R}\sqrt{A_j}}{\sqrt{2qI_d}}.
    \label{eq:dstar}
\end{align}

Generation-recombination noise from a trap ensemble with characteristic
lifetime $\tau_{\mathrm{gr}}$ has a Lorentzian form that falls below relevance
in the tens-of-GHz band of interest. Optical-source noise enters as:
\begin{equation}
    S_{\mathrm{RIN}}(f) = \mathcal{R}^2P_{\mathrm{opt}}^2\,\mathrm{RIN},
    \label{eq:rin_noise}
\end{equation}
and can become important at the link level even when the detector itself is quiet.
By contrast, taper insertion loss, waveguide propagation loss, free-carrier
absorption in the doped contact layer, Fresnel reflection, and incomplete carrier
collection do not create a new electrical noise PSD by themselves; they reduce
the signal term $I_{\mathrm{ph}}=\mathcal{R}P_{\mathrm{opt}}$ and thereby
degrade SNR indirectly.

\begin{figure}[H]
    \centering
    \missingfigure[figwidth=0.85\linewidth]{
        Two-panel photodetector-noise plot. Left panel: current-noise PSD
        showing $S_{\mathrm{shot}}=2qI_d$ and the load-defined thermal term
        $S_{\mathrm{thermal}}=4kT/R_L$ across frequency. Right panel: intrinsic
        photodetector noise variance $\sigma_n^2=i_n^2 B_e$ versus bandwidth,
        together with the corresponding dark-current shot-noise floor.}
    \caption{Noise decomposition of the photodetector from
    \cref{eq:shot,eq:thermal,eq:sigma_n}. For $I_d=\SI{1.3}{\nano\ampere}$,
    the dark-current shot-noise contribution remains small at
    \SI{103}{\giga\hertz}, confirming that the intrinsic detector noise floor
    is low.}
    \label{fig:noise_decomposition}
\end{figure}

\subsection[The Speed--Efficiency--Noise Trade-off]{The Speed--Efficiency--Noise Trade-off Triangle}
\label{subsec:pd_tradeoffs}

The metrics defined above are not free parameters. They are bound together by the
device physics into a multi-dimensional trade-off space that every design decision
navigates.

The fundamental tension in the responsivity-bandwidth trade-off is structural.
A thicker intrinsic layer $W_i$ absorbs more photons, improving $\mathcal{R}$
via \cref{eq:qe_full}, but also increases the carrier transit time, reducing
$f_{\mathrm{tt}}$ via \cref{eq:bw_tt}. A longer device absorbs more light but
increases $C_j$ linearly through \cref{eq:cj}, reducing $f_{\mathrm{RC}}$ via
\cref{eq:bw_rc}. Reducing $A_j$ below the absorption saturation point recovers
bandwidth at the cost of responsivity.

The noise axis adds a third dimension: higher reverse bias deepens the depletion
region, which raises $f_{\mathrm{tt}}$ by increasing the carrier saturation
velocity, but also drives more SRH generation through \cref{eq:jdark_srh},
degrading NEP via \cref{eq:nep}.

The responsivity-bandwidth product $\mathcal{R}\times f_{3\,\mathrm{dB}}$ captures
the most important two-dimensional trade-off as a scalar figure of merit and is
used in the benchmarking table of \cref{subsec:pd_benchmarks}. The design
rationale developed in \cref{sec:pd_design} maps the full
geometry-to-performance relationship, identifies the Pareto front of the
vertical architecture, and positions the chosen device within it.
